% =============================================================================
% Lecture 16: Deployment Release
% Physical AI: Machine Learning Systems That Sense and Act (Volume IV), Chapter 16
% =============================================================================
% Thesis: permission to operate is a claim about every condition the machine
% may meet next, so release states which physical claims the evidence
% supports, refuses wherever a premise is unknown, and expires when a
% premise's justification does.
% Running case: the warehouse mobile manipulator (the running machine) and its
% restricted mug-delivery case, claims C0 to C6. Real fleet: Waymo.
% Hook and synthesis: Challenger, 28 January 1986 (DR War Story 16.1).
%
% Build: make deck V=4 D=16_release   (figures: images/make_figs.py,
%        numbers: images/calcs.py)
% Plan:  lecture_map.yml in this folder.
% =============================================================================

\documentclass[aspectratio=169, 12pt]{beamer}
\usepackage{../../assets/beamerthememlsys}
\mlsyssetup{
  volume={Volume IV},
  chapter={Chapter 16},
  chaptertitle={Deployment Release},
  layer=governance,
  logo={../../assets/img/logo-mlsysbook.png},
  instlogo={../../assets/img/logo-harvard.png},
}
\usepackage[T1]{fontenc}
\usepackage[scaled=0.9]{helvet}
\usepackage{courier}
\renewcommand{\familydefault}{\sfdefault}
\usepackage{booktabs,amsmath}
\graphicspath{{images/}}
% Generated by calcs.py. Do not edit; edit calcs.py and rerun it.
\defval{air_f}{36}
\defval{colder_f}{15}
\defval{coldest_f}{53}
\defval{warmest_f}{81}
\defval{thiokol_f}{53}
\defval{n_flights}{23}
\defval{n_prior}{24}
\defval{flight_no}{25}
\defval{rings_words}{six}
\defval{breakup_s}{73}
\defval{joint_f}{28}
\defval{joint_err_f}{5}
\defval{cold_max_f}{63}
\defval{warm_min_f}{66}
\defval{cold_hit}{4}
\defval{cold_n}{4}
\defval{warm_hit}{3}
\defval{warm_n}{19}
\defval{n_distressed}{7}
\defval{t_lease}{60}
\defval{t_tick}{1}
\defval{t_bus}{1}
\defval{t_onset}{20}
\defval{t_silent}{82}
\defval{v_aisle}{1.3}
\defval{d_clear_m}{1.10}
\defval{d_clear_mm}{1{,}100}
\defval{f_latch}{100}
\defval{v_latch}{0.03}
\defval{f_contact}{50}
\defval{v_handover}{0.10}
\defval{wcet}{250}
\defval{deadline}{400}
\defval{payload}{50}
\defval{mu_min}{0.204}
\defval{tau_delay}{133.6}
\defval{a_brake}{2}
\defval{j_brake}{7.11}
\defval{allowances}{190}
\defval{d_loc}{50}
\defval{d_track}{40}
\defval{d_margin}{100}
\defval{reaction}{173.7}
\defval{scurve}{605.3}
\defval{budget}{969.0}
\defval{spare}{131.0}
\defval{t_spare}{100.8}
\defval{a_min}{1.39}
\defval{shortfall_pct}{31}
\defval{t_da}{63.3}
\defval{travel_da}{82.3}
\defval{reserve}{48.7}
\defval{t_reserve}{37.5}
\defval{wall_rate_exp}{-9}
\defval{confidence_pct}{95}
\defval{wall_h_mant}{3.0}
\defval{wall_h_exp}{9}
\defval{mu_oil}{0.05}
\defval{n_premises_words}{ten}
\defval{p_premise}{0.99}
\defval{p_conj}{0.904}
\defval{p_leg}{0.90}
\defval{p_leg_miss}{0.1}
\defval{p_leg_miss2}{0.10}
\defval{p_legs}{0.99}
\defval{t_exp}{16}
\defval{t_exp_dim}{32}
\defval{dt_age}{8}
\defval{dim_travel}{10.4}
\defval{dim_share_pct}{8}
\defval{scan_fast}{80}
\defval{scan_std}{160}
\defval{scan_fast_mm}{104}
\defval{scan_std_mm}{208}
\defval{f_trip}{15}
\defval{f_onset}{39}
\defval{x_rem}{0.15}
\defval{a_arm}{3}
\defval{t_arm}{10}
\defval{v_latch_hi}{0.10}
\defval{f_onset_hi}{95}
\defval{a_arm_hi}{400}
\defval{a_arm_ratio}{136}
\defval{a_arm_hi_g}{41}
\defval{g_one}{9.8}
\defval{f_handover}{20.7}
\defval{t_age}{51.6}
\defval{mu_insp}{0.12}
\defval{g}{9.81}
\defval{a_insp}{1.18}
\defval{ceiling_insp}{1.23}
\defval{v_restricted}{1}
\defval{waymo_cut}{June 2026}
\defval{waymo_miles}{271}
\defval{inj_n}{183}
\defval{inj_lo}{79}
\defval{inj_hi}{85}
\defval{ssi_n}{3}
\defval{ssi_lo}{83}
\defval{ssi_hi}{99}
\defval{fat_n}{2}
\defval{fat_lo}{18}
\defval{fat_hi}{98}
\defval{ssi_over_inj}{2.7}
\defval{mm_per_ms}{1.3}
\defval{d_silent}{106.6}
\defval{t_flagged}{22}
\defval{d_flagged}{28.6}
\defval{d_saved}{78}
\defval{t_crossover}{60}
\defval{rc1_n}{444}
\defval{rc2_n}{672}
\defval{rc3_n}{1{,}212}
\defval{rc_words}{three}
\defval{rc_months}{15}
\defval{rc_growth}{2.7}
\defval{rc_days}{137}
\defval{m_loaded}{368}
\defval{ke}{311}

\setcounter{mlsystotalsections}{5}

\title{Deployment Release}
\author{Vijay Janapa Reddi}
\institute{Harvard University}
\date{}

% Diagram role colors (Volume IV palette), for text beside figures
\definecolor{rolegreen}{HTML}{047857}
\definecolor{rolered}{HTML}{B91C1C}
\definecolor{roleamber}{HTML}{B45309}
\definecolor{roleblue}{HTML}{1D4ED8}
% A large answer choice for predict and concept-question slides
\newcommand{\choice}[2]{\textbf{\textcolor{mlsysaccent}{#1}}\enspace #2}
\newcommand{\answerat}[3]{\alt<#1->{\colorbox{mlsysaccent!15}{\choice{#2}{#3}}}{\choice{#2}{#3}}}
% A goal or takeaway line with its text-free thumbnail
\newcommand{\thumbline}[2]{%
  \begin{minipage}[c]{3.3cm}\centering\includegraphics[width=3.2cm,height=1.15cm,keepaspectratio]{#1}\end{minipage}%
  \hspace{0.3cm}\begin{minipage}[c]{\dimexpr\linewidth-3.7cm\relax}\small #2\end{minipage}}
% One line of a worked example, shown from overlay #1 on
\newcommand{\wstep}[2]{\uncover<#1->{#2}\\[0.15cm]}
% An exact excerpt of a record, datasheet, or metadata file
\newenvironment{logblock}[1][7]{\par\ttfamily\fontsize{#1}{8.4}\selectfont\raggedright\setlength{\parskip}{0pt}}{\par}
\newcommand{\ind}[1]{\hspace*{#1\fontcharwd\font`0}}

\begin{document}

% =============================================================================
% TITLE AND HOOK
% =============================================================================
\mlsysnote{[0.5 min] [CORE] Read the thesis aloud once. This is the third lecture of Part IV, whose law is the fourth law, permission must be earned. Lecture 15 left the machine with signed fault records; today decides what those records, and every earlier record, license. Notes total about 71 minutes; the core is about 60. FLEX: notes tagged [OPTIONAL] can be cut; cut first the Thiokol documents, the update-metadata excerpt, and the erosion history. The Waymo intervals are core: they carry the Expiry segment's real data.}
\mlsystitle{Deployment Release}{Release states which physical claims the evidence supports, refuses wherever a premise is unknown, and expires when a premise's justification does.}{cover_release_blueprint.png}

\begin{frame}{Predict: what did \val{n_flights} flights say about \val{air_f}~°F?}
  \mlsysactivity{Predict}{2 min}
  \begin{columns}[c]
    \begin{column}{0.30\textwidth}
      \centering\safeimg[width=\textwidth,height=4.4cm,keepaspectratio]{challenger-icicles.jpg}\\[0.05cm]
      {\footnotesize Launch tower, 28 January 1986}
    \end{column}
    \begin{column}{0.66\textwidth}
      {\small Air at launch: \textbf{\val{air_f}~°F}. Coldest joint of \val{n_flights} inspected flights: \textbf{\val{coldest_f}~°F}.\\
      The motor maker's engineers had recommended no launch below \val{thiokol_f}~°F.}\\[0.15cm]
      \choice{A}{Safe: every earlier flight came back}\\[0.05cm]
      \choice{B}{Unsafe: cold flights damaged seals}\\[0.05cm]
      \answerat{2}{C}{Nothing: no flight was that cold}\\[0.1cm]
      \uncover<2->{\begin{columns}[c]
        \begin{column}{0.42\textwidth}\safeimg[width=\textwidth,height=1.7cm,keepaspectratio]{challenger-srb-smoke.jpg}\end{column}
        \begin{column}{0.56\textwidth}\footnotesize Smoke at the right booster's aft field joint at ignition. The vehicle broke up about \val{breakup_s}~s later.\end{column}
      \end{columns}}
    \end{column}
  \end{columns}
  \mlsyscredit{Photos: NASA, public domain, Wikimedia Commons. Deployment Release, War Story 16.1; Rogers Commission report, 1986, ch.~IV finding 6.}
  \note{[3 min] [CORE] NARRATE: Space Shuttle Challenger, 28 January 1986. Icicles hung on the launch tower. The air at launch was \val{air_f} degrees F, \val{colder_f} degrees colder than any earlier launch (Rogers Commission). Challenger was the \val{flight_no}th flight; of the \val{n_prior} before it, \val{n_flights} had their boosters recovered and inspected, so the seal record holds \val{n_flights} flights. The booster's field joints were sealed by two rubber O-rings; the coldest earlier flight had joints at \val{coldest_f} F and the worst seal damage to date. On the eve of launch, engineers at Morton Thiokol, the motor maker, recommended against launching below \val{thiokol_f} F. ENGAGE: vote A, B, or C. Step 2: C is the release answer. The record contains no flight at that temperature, so the premise the launch needed was unknown, and an unknown premise refuses. B is what the whole record suggests, as the next segment shows, but C is the reasoning this lecture builds. Students who pick A treat a success count as evidence about conditions it never sampled. Point to the smoke photo: a gas leak at the right aft joint at ignition; all seven crew died. WARN: no learned model is in this story; the failure is an approval for conditions the evidence did not cover. FLEX: [CORE]; the synthesis returns here.}
\end{frame}

\begin{frame}{Recall: what does a fault record prove?}
  \mlsysactivity{Retrieval}{1.5 min}
  \begin{columns}[c]
    \begin{column}{0.56\textwidth}
      {\Large\bfseries Question of the day}\\[0.2cm]
      {\normalsize What does one fault record contain? When the proposer stalls, how soon must the base brake?}\\[0.3cm]
      \uncover<2->{\small An injected fault, an oracle fixed before the run, and the observed trace. The set also names the regimes no trial reached.\\[0.1cm]
      Proposer stall: brake onset within \textbf{\val{t_silent}~ms} of the last valid renewal.}
    \end{column}
    \begin{column}{0.40\textwidth}
      \centering\safeimg[width=\textwidth,height=1.6cm,keepaspectratio]{thumb_silent_step2.pdf}\\[0.05cm]
      {\footnotesize Lease \val{t_lease} + tick \val{t_tick} + bus \val{t_bus} + onset \val{t_onset} = \val{t_silent}~ms}\\[0.2cm]
      \uncover<3->{\begin{mlsyscard}{mlsysaccent}\small\textbf{Today:} which claims do these records license, under which conditions, and for how long?\end{mlsyscard}}
    \end{column}
  \end{columns}
  \mlsyscredit{Lecture 15, Adversarial Verification; Deployment Release, \S16.1 and Napkin Math 16.2.}
  \note{[1.5 min] [CORE] NARRATE: give 45 seconds. Step 2: Lecture 15's fault record pairs a trace under an injected fault (a stalled proposer, a stale frame, lost fieldbus frames, a sagging bus) with an oracle fixed before the trial. The set also lists the regimes no trial reached. For fault F1, the stalled proposer, brake onset had to come within \val{t_silent} ms: the \val{t_lease} ms chunk lease runs out, one \val{t_tick} ms tick and one \val{t_bus} ms bus transfer pass, and the brake takes \val{t_onset} ms to engage. Step 3: a fault record proves nothing outside the envelope it sampled, and an untested mode leaves a gap in coverage. ENGAGE: is a missing fault record a PASS? (Expected: no; it leaves the claim unknown.) FLEX: [CORE].}
\end{frame}

\begin{frame}{Why can a machine that passes every trial still be refused?}
  \thumbline{thumb_budget_bar_aisle_step8.pdf}{Write release as physical inequalities, and find how much clear distance the aisle stop leaves.}\\[0.18cm]
  \thumbline{thumb_legs_step2.pdf}{Build a safety case, and find the premise that limits it.}\\[0.18cm]
  \thumbline{thumb_tree_step3.pdf}{Sign UNCOND, COND, or REFUSE, and name the measurement that would change it.}\\[0.18cm]
  \thumbline{thumb_manifest_step3.pdf}{Bind the verdict to one configuration and its standing conditions, and say what an update reopens.}\\[0.15cm]
  {\footnotesize\textcolor{midgray}{Thread: the fourth law, permission must be earned. Record: the release manifest.}}
  \note{[1.5 min] [CORE] The driving question: a machine can complete every trial and still lack the evidence to be allowed among people. Four goals, one per part of the lecture. The fourth law, permission must be earned, which Lecture 14 stated, is a design principle the book elevates to a bedrock law; only the first law, irreversibility, is physical. Today adds one record to Lecture 1's record chain, the release manifest, the record that binds every earlier record to one configuration. The running case is the warehouse mobile manipulator delivering a mug; the real fleet beside it is Waymo. FLEX: [CORE].}
\end{frame}

% =============================================================================
% SECTION: CLAIMS
% =============================================================================
\mlsysnote{[0.5 min] [CORE] First, what a release decides, and why a record of successes cannot decide it.}
\mlsysdivider{What does a release decide, and on what evidence?}{Deployment Release, \S16.1 to \S16.3.1}
\section{Claims}

\begin{frame}{Every flight in the record had joints at \val{coldest_f}~°F or warmer}
  \centering
  \mlsyssteps[width=0.97\textwidth,height=5.0cm,keepaspectratio]{oring_step}{3}
  \mlsyscredit{Data: Dalal, Fowlkes, and Hoadley, JASA 1989, table 1, via the UCI Machine Learning Repository; temperatures as in Rogers Commission figure 7. Challenger: Rogers ch.~IV finding 6 (joint \val{joint_f} $\pm$ \val{joint_err_f}~°F).}
  \note{[2.5 min] [CORE] Real data: \val{n_flights} flights before Challenger, each with its calculated joint temperature at launch and the number of its \val{rings_words} O-rings that showed erosion or blow-by. Step 1: only the flights with distress, the view reviewers used. Temperature shows no clear pattern; the Commission's figure 6 drew exactly this. Step 2: all \val{n_flights} flights. Every flight at \val{cold_max_f} F or colder had distress (\val{cold_hit} of \val{cold_n}); at \val{warm_min_f} F or warmer, \val{warm_hit} of \val{warm_n} (counted in calcs.py). Step 3: the shaded region holds no flight at all. Challenger's coldest joint point was estimated at \val{joint_f} F, plus or minus \val{joint_err_f}, with the air at \val{air_f} F. ENGAGE: what claim does this record support at \val{joint_f} F? (Expected: none; the premise is unknown.) WARN: students will want to fit a curve and extrapolate. A fit is an argument, and its warrant fails outside the data. LINK: this is the envelope idea, in which a claim covers only the conditions its evidence sampled. FLEX: [CORE].}
\end{frame}

\begin{frame}{Release decides three physical inequalities for one machine}
  \begin{columns}[c]
    \begin{column}{0.36\textwidth}
      \centering\safeimg[width=\textwidth,height=4.3cm,keepaspectratio]{erflex-mobile-manipulator.jpg}\\[0.05cm]
      {\footnotesize A real machine of the class: MiR250 base, UR16e arm}
    \end{column}
    \begin{column}{0.62\textwidth}
      \small
      \textbf{Stop:} $d_{\text{stop}} \le D_{\text{clear}}$, base at \textbf{\val{v_aisle}~m/s}, clear distance \textbf{\val{d_clear_m}~m}\\[0.15cm]
      \textbf{Contact:} latch force $\le$ \textbf{\val{f_latch}~N} at \val{v_latch}~m/s; handover $\le$ \textbf{\val{f_contact}~N} at \val{v_handover}~m/s\\[0.15cm]
      \textbf{Timing:} permission cycle WCET $\le$ deadline, \textbf{\val{wcet} of \val{deadline}~$\mu$s}\\[0.25cm]
      \uncover<2->{\footnotesize Envelope: payload $\le$ \val{payload}~kg, floor friction $\mu \ge$ \val{mu_min}, surveyed rack ends, person standing still, local authority only.\\[0.1cm]
      WCET: worst-case execution time.}
    \end{column}
  \end{columns}
  \mlsyscredit{Deployment Release, equation 1 and table 1 (speeds chosen; limits illustrative). Photo: Auledas, CC BY 4.0, Wikimedia Commons (ER-FLEX by Enabled Robotics).}
  \note{[2.5 min] [CORE] LINK: the running machine of every lecture: a wheeled base with an articulated arm. The photo shows a real machine of that class, the ER-FLEX, a MiR250 mobile base carrying a UR16e arm; our machine is a reference platform with its own numbers. NARRATE: release turns into three testable inequalities. The base stops inside the clear distance at the \val{v_aisle} m/s aisle speed. Contact stays under its limit: \val{f_latch} N on the cage-door latch at a guarded \val{v_latch} m/s approach, and \val{f_contact} N on a coworker's hand at a \val{v_handover} m/s handover; the \val{f_contact} N is a site choice; no injury threshold sets it. The permission path finishes each check inside its deadline: a declared \val{wcet} of \val{deadline} microseconds. Step 2: the operating envelope that every claim assumes. ENGAGE: which envelope item could change between Monday and Tuesday without any software change? (Expected: floor friction, payload, a person who keeps walking.) WARN: a benchmark score or a demo video decides none of these inequalities. FLEX: [CORE].}
\end{frame}

\begin{frame}{Estimate: how much clear distance is left?}
  \mlsysactivity{Estimate}{2 min}
  \centering
  \only<1-2|handout:0>{\mlsyskitfig[width=0.95\textwidth,height=4.8cm,keepaspectratio]{budget_bar_aisle_step7}}\only<3-|handout:1>{\mlsyskitfig[width=0.95\textwidth,height=4.8cm,keepaspectratio]{budget_bar_aisle_step8}}\\[0.05cm]
  \footnotesize
  \wstep{2}{Spare \val{spare}~mm $=$ \val{t_spare}~ms of extra delay, or braking at \textbf{\val{a_min}}~m/s$^2$ (\val{shortfall_pct}\% weaker)}
  \uncover<3->{Admission work, \val{v_aisle}~m/s $\times$ \val{t_da}~ms $=$ \val{travel_da}~mm, leaves a reserve of \textbf{\textcolor{rolered}{\val{reserve}~mm}}}
  \mlsyscredit{Deployment Release, Napkin Math 16.1 (values illustrative or chosen; derived in calcs.py). Dispatch-to-admission: Lecture 13. Stopping-budget bar: course kit, MLSysIM values.}
  \note{[3.5 min] [CORE] LINK: Lecture 12 built this budget and chose \val{v_aisle} m/s; release inherits its terms as premises. ENGAGE: the bar shows \val{spare} mm of the \val{d_clear_mm} mm left. Two minutes in pairs. What single wrong input would spend it, in milliseconds of delay or in deceleration? Step 1: the inputs are Lecture 12's, \val{v_aisle} m/s, \val{tau_delay} ms of delay, \val{a_brake} m/s squared, \val{j_brake} m/s cubed of jerk, and \val{allowances} mm of allowances. Reaction travel \val{v_aisle} m/s times \val{tau_delay} ms = \val{reaction} mm; the jerk-limited stop (Lecture 11) brakes over \val{scurve} mm. Step 2: \val{d_loc} mm localization, \val{d_track} mm tracking, \val{d_margin} mm margin: \val{budget} mm used, \val{spare} mm spare. One input wrong: \val{spare} mm is \val{t_spare} ms of extra delay, or a deceleration of \val{a_min} instead of \val{a_brake} m/s squared, \val{shortfall_pct} percent weaker (derived). Step 3: Lecture 13's dispatch-to-admission allocation, \val{t_da} ms, costs \val{travel_da} mm; the reserve is \val{reserve} mm, \val{t_reserve} ms at aisle speed. WARN: the margin that matters is the \val{spare} mm left unspent. A brake lining \val{shortfall_pct} percent weaker spends it all. FLEX: [CORE]. This is worked example 1.}
\end{frame}

\begin{frame}{A safety case links each claim to evidence through a warrant}
  \centering
  \mlsyssteps[width=0.97\textwidth,height=5.5cm,keepaspectratio]{cae_step}{3}
  \mlsyscredit{After Deployment Release, \S16.3.1 to \S16.3.3 (claim-argument-evidence; Goal Structuring Notation; Toulmin 1958). Values illustrative.}
  \note{[2.5 min] [CORE] Definition: a safety case is a structured argument that a named machine, in a stated envelope and fault model, meets explicit physical criteria. Claim-argument-evidence (CAE) and Goal Structuring Notation (GSN) draw the same tree. Step 1: the claim at the top, the context (the envelope) beside it, the warrant (why the evidence supports the claim: the stopping budget bounds travel if every term is bounded), and evidence leaves: calibrated records with stated uncertainty. Step 2: 1,000 clean runs are a demonstration. They record what happened, but no warrant carries them to conditions they did not sample. LINK: Lecture 7's exposure wall: showing 1e\val{wall_rate_exp} failures per hour with zero failures at \val{confidence_pct} percent takes \val{wall_h_mant}e\val{wall_h_exp} hours (derived). So the case rests on an architectural warrant, the independent permission path; trial counts cannot carry it. Step 3: a defeater is a real condition that breaks the warrant, such as a clear oil film near \val{mu_oil} friction. ENGAGE: name another defeater for this warrant. (Expected: a person who keeps walking; a worn brake.) FLEX: [CORE]. Land the segment: a claim is an inequality over an envelope, and the case says why the evidence supports it.}
\end{frame}

% =============================================================================
% SECTION: PREMISES
% =============================================================================
\mlsysnote{[0.5 min] [CORE] A case has many premises. How strong is it, and who decides which standards it must meet?}
\mlsysdivider{Why is a case no stronger than its weakest premise?}{Deployment Release, \S16.3.2 to \S16.3.4}
\section{Premises}

\begin{frame}{Required premises multiply; only independent legs add}
  \begin{columns}[c]
    \begin{column}{0.64\textwidth}
      \centering\mlsyssteps[width=\textwidth,height=5.5cm,keepaspectratio]{legs_step}{2}
    \end{column}
    \begin{column}{0.34\textwidth}
      \centering\safeimg[width=\textwidth,height=3.3cm,keepaspectratio]{srb-field-joint.jpg}\\[0.05cm]
      {\footnotesize Booster field joint: primary and secondary O-ring. In 1982 NASA reclassified the pair from redundant (1R) to single-point (1): joint rotation could unseat both.}
    \end{column}
  \end{columns}
  \mlsyscredit{Deployment Release, equations 2 and 3 (stipulated model). Drawing: Rogers Commission, vol.~1, p.~57, public domain; reclassification: Rogers ch.~VI.}
  \note{[3 min] [CORE] Step 1: when every child claim is required, the case is conjunctive. In a stipulated model where each premise holds with probability \val{p_premise}, \val{n_premises_words} premises hold together at \val{p_conj}, below every single one (derived). The stopping claim has this form, since clear distance, deceleration, brake onset, observation age, localization, and tracking all enter one budget. An assessor therefore looks for the least supported premise; the volume of documents settles nothing. Step 2: two sufficient legs, either of which proves the claim, give 1 minus \val{p_leg_miss} squared, \val{p_legs}, but only if they fail independently. Point to the joint: the primary and secondary O-rings looked like two legs, and the Commission found that once the joint rotated under pressure, both could fail; Marshall changed the criticality from 1R to 1 in December 1982. WARN: policy validation and a permission path fail together while they share processors, memory, or power rails (claim C5; Lecture 13's resource inventory). Knight and Leveson (1986) measured coincident failures in independently written programs. FLEX: [CORE].}
\end{frame}

\begin{frame}{Vote: which statements hold?}
  \mlsysactivity{Vote}{2 min}
  \normalsize True or false?\\[0.3cm]
  \choice{1}{\Val{n_premises_words} premises at \val{p_premise} each support the claim at \val{p_premise}.}\uncover<2->{\enspace\textbf{\textcolor{rolered}{False}}}\\[0.25cm]
  \choice{2}{A validated policy and an enforcer on the same chip are two independent legs.}\uncover<2->{\enspace\textbf{\textcolor{rolered}{False}}}\\[0.25cm]
  \choice{3}{High detection accuracy can offset an unmeasured braking premise.}\uncover<2->{\enspace\textbf{\textcolor{rolered}{False}}}\\[0.25cm]
  \choice{4}{Dim light that ages frames \val{dt_age}~ms is a premise the case must state.}\uncover<2->{\enspace\textbf{\textcolor{rolegreen}{True}}}
  \mlsyscredit{Deployment Release, \S16.3.2 (dim-aisle example: \val{t_exp} to \val{t_exp_dim}~ms exposure, \val{dim_travel}~mm at \val{v_aisle}~m/s).}
  \note{[2 min] [CORE] Peer instruction: vote, one minute with a neighbor, vote again. Step 2: (1) false; \val{p_conj}. (2) false; shared memory buses, caches, and rails let one fault reach both, so independence must be demonstrated. (3) false; in a safety case one branch cannot compensate for another. (4) true; if dim light doubles camera exposure from \val{t_exp} to \val{t_exp_dim} ms, frames arrive \val{dt_age} ms older, \val{dim_travel} mm of travel at \val{v_aisle} m/s, about \val{dim_share_pct} percent of the \val{spare} mm spare. On this machine the observation contract's freshness check (Lecture 8) refuses the old frames, so the cost shows up as lost availability, with no collision. A camera mount that drifts while frames stay fresh is caught by nothing. FLEX: [CORE].}
\end{frame}

\begin{frame}{Certification authorizes one configuration with its evidence}
  \begin{columns}[T]
    \begin{column}{0.50\textwidth}
      \footnotesize
      \begin{tabular}{@{}P{2.1cm}P{4.8cm}@{}}\toprule
        \textbf{Standard} & \textbf{Where it applies on the machine}\\\midrule
        ISO 13849-1 & Safety MCU's lease watchdog: the one rated function in the permission path\\
        ISO 3691-4 & The base as a driverless industrial truck\\
        ISO 10218-2 & The arm; excludes the arm moving on the base\\
        ANSI/A3 R15.08 & The whole machine (mobile robot with a manipulator)\\\bottomrule
      \end{tabular}
    \end{column}
    \begin{column}{0.48\textwidth}
      {\small\textbf{Keyence SZ-V32 safety laser scanner}}\par\vspace{0.1cm}
      \begin{logblock}[8.6]
        Type 3; PL d, Category 3\\[0.1cm]
        Protection zone (70/150 mm objects):\\
        \ind{2}8.4 m standard /\\
        \ind{2}5.7 m high-speed mode\\[0.1cm]
        Response time (ON to OFF):\\
        \ind{2}160-1280 ms standard\\
        \ind{2}\textbf{80-640 ms high-speed}\\
        \ind{2}(2 to 16 scans)
      \end{logblock}
      \vspace{0.1cm}
      \uncover<2->{\footnotesize At \val{v_aisle}~m/s: \val{scan_fast}~ms is \textbf{\val{scan_fast_mm}~mm} of travel; \val{scan_std}~ms is \textbf{\val{scan_std_mm}~mm} (derived).}
    \end{column}
  \end{columns}
  \mlsyscredit{Deployment Release, table 3 and \S16.3.4. Datasheet: Keyence SZ-V series specifications, as read into the MLSysIM registry (2026-09-29).}
  \note{[2.5 min] [CORE] NARRATE: this machine is industrial machinery, so certification follows the machinery route. ISO 12100 sets the risk assessment; ISO 13849-1 rates control functions by performance level; ISO 3691-4 covers the base; ISO 10218-2 covers the arm but excludes a manipulator riding on a mobile platform; ANSI/A3 R15.08 treats the combination as one industrial mobile robot. Within the permission path only the safety MCU's lease and heartbeat watchdog is rated; the admission checks are not. Right: a real scanner datasheet. Its rating and its response time hold for a configured mode, object size, and scan count. Step 2: the same scanner in standard mode takes twice as long, \val{scan_std_mm} against \val{scan_fast_mm} mm of travel at aisle speed. ENGAGE: the integrator switches the scanner to standard mode for range. What happens to the certificate? (Expected: it no longer covers that configuration; the stopping claim reopens.) WARN: under the EU Machinery Regulation an unforeseen update that raises risk is a substantial modification and can make the updater the manufacturer (\S16.3.4). FLEX: [CORE]. Land the segment.}
\end{frame}

% =============================================================================
% SECTION: VERDICTS
% =============================================================================
\mlsysnote{[0.5 min] [CORE] The argument must end in a decision the machine can enforce. What can the adjudicator sign?}
\mlsysdivider{What may the adjudicator sign when a premise is unknown?}{Deployment Release, \S16.4 and \S16.6}
\section{Verdicts}

\begin{frame}{One evening, the same data, two signed recommendations}
  \begin{columns}[c]
    \begin{column}{0.47\textwidth}
      \centering\safeimg[width=\textwidth,height=2.6cm,keepaspectratio]{thiokol-recommendation.jpg}\\[0.05cm]
      {\footnotesize Thiokol engineering chart, 27 January: ``O-ring temp must be $\ge$ \val{thiokol_f}~°F at launch}\\[0.25cm]
      \uncover<2->{\footnotesize\itshape ``We had to prove to them that we weren't ready.''}\\
      \uncover<2->{\footnotesize Robert Lund, Thiokol vice president of engineering, testimony}
    \end{column}
    \begin{column}{0.51\textwidth}
      \centering\uncover<2->{\safeimg[width=\textwidth,height=4.7cm,keepaspectratio]{thiokol-launch-fax.jpg}\\[0.02cm]
      {\footnotesize Later that night: ``MTI recommends STS-51L launch proceed}}
    \end{column}
  \end{columns}
  \mlsyscredit{Rogers Commission report, vol.~1, ch.~V: p.~90 (chart), p.~97 (assessment), Lund testimony. Public domain.}
  \note{[2.5 min] [OPTIONAL] Two real documents. Step 1: Thiokol's engineers ended their chart with a recommendation. O-ring temperature must be at least \val{thiokol_f} F at launch, because the closest evidence, flight SRM 15, worked at \val{coldest_f} F. That is a REFUSE that names its premise. Marshall managers rejected the rationale. Step 2: when the teleconference resumed at about 11 p.m., Thiokol management, having reassessed, signed an assessment stating that the temperature data are not conclusive, so MTI recommends launch proceed on 28 January. ENGAGE: what changed between the two documents? (Expected: no data; only who carried the burden of proof.) Lund testified: we had to prove to them that we weren't ready. The burden had inverted. The reviewers demanded proof of danger outside the evidence, where no data existed, in place of proof of safety inside it. WARN: the fix is structural. A verdict keeps unknown as unknown, and unknown refuses. FLEX: [OPTIONAL]; the synthesis returns to this point.}
\end{frame}

\begin{frame}{A verdict is a signed permit for a stated envelope}
  \begin{columns}[c]
    \begin{column}{0.42\textwidth}
      \begin{mlsyscard}{rolegreen}\footnotesize\textbf{UNCOND:} the full assessed envelope\end{mlsyscard}
      \vspace{0.02cm}
      \begin{mlsyscard}{roleamber}\footnotesize\textbf{COND:} a restricted envelope with its own evidence\end{mlsyscard}
      \vspace{0.02cm}
      \begin{mlsyscard}{rolered}\footnotesize\textbf{REFUSE:} a premise missing, contradicted, or unknown; no permit\end{mlsyscard}
      \vspace{0.05cm}
      {\footnotesize Each lasts while its standing conditions hold. No verdict moves the proposal boundary.}
    \end{column}
    \begin{column}{0.56\textwidth}
      \centering\safeimg[width=0.80\textwidth,height=2.6cm,keepaspectratio]{drive-pilot-roads.jpg}\\[0.05cm]
      {\footnotesize DRIVE PILOT approved roads, California and Nevada, June 2024}\\[0.08cm]
      \uncover<2->{\begin{logblock}[6.6]
        ``can only operate on highways during daylight at speeds not exceeding 40 miles per hour. This permit excludes [...] construction zones, during heavy rain or heavy fog [...]''
      \end{logblock}}
    \end{column}
  \end{columns}
  \mlsyscredit{Deployment Release, \S16.4 and table 4. Map: Tamnguyenrn, CC BY-SA 4.0, Wikimedia Commons. Text: California DMV press release, 8 June 2023.}
  \note{[2.5 min] [CORE] Definition: a release verdict is the signed decision that sets how much authority the permission path may grant. UNCOND covers the full assessed envelope; COND covers a separately evidenced restricted envelope; REFUSE grants no permit and commands the assessed stop or hold. In this book's architecture the verdict is a signed descriptor loaded into the safety microcontroller, and every verdict lasts only while its standing conditions hold. Under UNCOND and COND the learned policy still only proposes; the enforcer still admits each command. Right: a real COND. In June 2023 the California DMV permitted Mercedes-Benz DRIVE PILOT, a Level 3 system, under a restricted envelope of highways, daylight, at most 40 mph, no construction zones, rain, or fog. The map shows the approved roads a year later. ENGAGE: is this permit a statement that the car drives well? (Expected: no; it names a configuration and the conditions its evidence covers.) WARN: COND keeps every premise; the restricted envelope needs its own evidence. FLEX: [CORE].}
\end{frame}

\begin{frame}{Predict: can the door branch be signed?}
  \mlsysactivity{Predict}{2 min}
  \begin{columns}[c]
    \begin{column}{0.55\textwidth}
      \centering\mlsyssteps[width=\textwidth,height=3.7cm,keepaspectratio]{door_step}{2}
    \end{column}
    \begin{column}{0.43\textwidth}
      \small Guarded approach: \textbf{\val{f_onset}~N} at brake onset; the latch allows \textbf{\val{x_rem}~mm} more before \val{f_latch}~N.\\[0.2cm]
      \choice{A}{UNCOND: \val{f_onset} is under \val{f_latch}}\\[0.08cm]
      \choice{B}{COND at \val{v_latch}~m/s only}\\[0.08cm]
      \answerat{3}{C}{REFUSE: arm stop unmeasured}
    \end{column}
  \end{columns}
  \vspace{0.1cm}
  \centering\uncover<3->{\footnotesize The arithmetic is exact; the premise is not. No trace shows the loaded arm stopping in \val{x_rem}~mm.}
  \mlsyscredit{Deployment Release, \S16.6 (Lecture 12's equations 12.11 and 12.12; values illustrative, computed with MLSysIM).}
  \note{[3 min] [CORE] LINK: Lecture 12 derived the contact bounds. Step 1: at the guarded \val{v_latch} m/s approach, the tripwire fires at \val{f_trip} N and the force reaches \val{f_onset} N by brake onset; the latch has \val{x_rem} mm left before \val{f_latch} N, so the arm must decelerate at about \val{a_arm} m/s squared and stop in about \val{t_arm} ms. ENGAGE: vote before step 2. Step 2: at \val{v_latch_hi} m/s the force at onset is \val{f_onset_hi} N and the stop needs \val{a_arm_hi} m/s squared, \val{a_arm_ratio} times more, about \val{a_arm_hi_g} g (derived). That is why the envelope admits only the guarded approach. Step 3: C. Even at \val{v_latch} m/s nobody has measured whether the arm, with its link momentum, payload, and reflected rotor inertia, can deliver \val{a_arm} m/s squared within a fraction of a millimeter. The premise is unknown, so the claim refuses. Students who pick A read the force check as the whole claim; B would need its own arm-stopping evidence. WARN: neither number is in doubt as arithmetic; refusal comes from the premise beneath the equation. FLEX: [CORE].}
\end{frame}

\begin{frame}{On the book's evidence every branch refuses}
  \centering
  \mlsyssteps[width=0.98\textwidth,height=5.7cm,keepaspectratio]{tree_step}{3}
  \mlsyscredit{After Deployment Release, figure 1 and table 2. Leaf values illustrative, declared, or unknown; none measured on the loaded configuration.}
  \note{[2.5 min] [CORE] The whole case for restricted mug delivery. Step 1: C0 holds only if all of C1 to C6 hold; C1 splits into base stop, door latch, and mug handover; each claim ends in its key premise. Step 2: the outline marks each premise's evidence: dashed illustrative, dash-dot declared (the \val{wcet} microsecond WCET has no analysis behind it), dotted unknown (no placement test report, no PASS records for the authority handovers A1 to A4 of Lecture 14). The oil film attacks C3. Step 3: every claim refuses, so the machine's verdict is REFUSE. The verdict is illustrative too, because the book's inputs are models; no test logs exist. ENGAGE: which single measurement would you buy first? (Expected: loaded stopping trials; they settle the deceleration behind C1 and C3.) Each refusal names its next measurement: loaded stops for C1, brake-onset trials for C2, an end-to-end age trace for C4, loaded enforcer timing for C5. C3's open premise passes to the residual-claims register of Lecture 18. FLEX: [CORE].}
\end{frame}

\begin{frame}{With stipulated evidence, mug delivery earns COND at \val{v_restricted}~m/s}
  \begin{columns}[c]
    \begin{column}{0.40\textwidth}
      \centering\safeimg[width=\textwidth,height=3.0cm,keepaspectratio]{robot-cell-fence.jpg}\\[0.05cm]
      {\footnotesize Fencing: one way a restricted case excludes contact}
    \end{column}
    \begin{column}{0.58\textwidth}
      \small
      Inspected floors, $\mu \ge \val{mu_insp}$: braking $\le$ \val{mu_insp} $\times$ \val{g} $=$ \textbf{\val{a_insp}~m/s$^2$}\\[0.08cm]
      Speed ceiling \textbf{\val{ceiling_insp}~m/s}; signed speed \textbf{\val{v_restricted}~m/s}\\[0.08cm]
      Local authority only; consenting coworker; singulated feeder\\[0.15cm]
      \uncover<2->{\footnotesize Any spill, expired inspection, or perimeter breach returns the verdict to \textbf{\textcolor{rolered}{REFUSE}}. A lower speed cannot repair an oil film near $\mu$ = \val{mu_oil}.}
    \end{column}
  \end{columns}
  \mlsyscredit{Deployment Release, \S16.4.1 and \S16.6.1 (stipulated evidence; ceilings derived in calcs.py). Photo: PEO ACWA, CC BY 2.0, Wikimedia Commons.}
  \note{[2.5 min] [CORE] Worked example 3, under stipulated evidence: suppose loaded stopping traces, latch and handover force trials, loaded timing and fault-injection logs, an end-to-end age trace, and PASS records for every local handover all arrive. Step 1: the floor regime certifies friction of only \val{mu_insp}, so the optimistic traction ceiling is \val{mu_insp} times \val{g} = \val{a_insp} m/s squared, and the same budget admits at most \val{ceiling_insp} m/s (derived). The case signs \val{v_restricted} m/s, below the ceiling, and charges dispatch-to-admission delay. The task packet must also pass predeclared success, cycle-time, and intervention criteria. The photo shows another restriction: a fence keeps people out of the arm's volume, if interlock and stopping evidence support it. Step 2: COND lives only while its premises hold. ENGAGE: why can't the site keep running at 0.5 m/s on a reported spill? (Expected: the speed was sized for \val{mu_insp}; a film removes most traction.) FLEX: [CORE]. Land the segment: unknown refuses; refusal names the measurement; COND needs its own evidence.}
\end{frame}

% =============================================================================
% SECTION: EXPIRY
% =============================================================================
\mlsysnote{[0.5 min] [CORE] A verdict is signed against the machine as it is today. Brakes wear, mounts shift, floors polish. When does a granted verdict lapse?}
\mlsysdivider{When does a granted verdict lapse?}{Deployment Release, \S16.5, \S16.7, \S16.9.5 and \S16.9.6}
\section{Expiry}

\begin{frame}{Predict: which crash outcome is least certain?}
  \mlsysactivity{Predict}{1.5 min}
  \begin{columns}[c]
    \begin{column}{0.70\textwidth}
      \centering\mlsyssteps[width=\textwidth,height=4.6cm,keepaspectratio]{waymo_step}{2}
    \end{column}
    \begin{column}{0.28\textwidth}
      \centering\safeimg[width=\textwidth,height=2.6cm,keepaspectratio]{waymo-jaguar-sf.jpg}\\[0.05cm]
      {\footnotesize Driverless Waymo, 2023}\\[0.15cm]
      {\footnotesize \val{waymo_miles} million miles bound common crashes tightly and fatal ones loosely}
    \end{column}
  \end{columns}
  \mlsyscredit{Waymo Safety Impact Data Hub, \val{waymo_cut} release (rider-only, location-adjusted benchmark), via MLSysIM; Kusano et al.\ 2025. Photo: Dllu, CC BY 4.0, Wikimedia Commons.}
  \note{[3 min] [CORE] Real fleet evidence, owned by this lecture (Lecture 18 recalls it in one line). Waymo publishes crash rates by outcome against a human benchmark for the same places, with 95 percent intervals, and republishes as miles accumulate. Step 1: common outcomes after \val{waymo_miles} million rider-only miles: police-reported crashes, injuries (\val{inj_n} crashes, a reduction of \val{inj_lo} to \val{inj_hi} percent), airbag deployments. ENGAGE: predict which of the rare outcomes has the widest interval. Step 2: serious injury rests on \val{ssi_n} crashes (\val{ssi_lo} to \val{ssi_hi} percent) and fatal on \val{fat_n} (\val{fat_lo} to \val{fat_hi} percent). The serious-injury interval is \val{ssi_over_inj} times the injury interval (derived). Rare events accumulate slowly however many miles a fleet drives. LINK: Lecture 7's exposure wall, measured on a real fleet. WARN: each outcome's claim may extend only as far as its own interval; read the interval before the point estimate, and never pool routes or count interventions as collisions. FLEX: [CORE]; it carries this segment's real data.}
\end{frame}

\begin{frame}{Each premise needs a detector, a cadence, and a response}
  \footnotesize
  \begin{tabular}{@{}P{2.6cm}P{3.9cm}P{3.1cm}P{3.6cm}@{}}\toprule
    \textbf{Premise} & \textbf{Detector} & \textbf{Cadence} & \textbf{When it fails}\\\midrule
    Speed, lease, tracking & encoders, IMU, lease timer & every \val{t_tick}~ms tick & clamp, or fallback stop\\
    Payload $\le$ \val{payload}~kg & rack load cells & each pick & refuse aisle transit\\
    Floor $\mu \ge$ \val{mu_min} & survey and spill procedure & site schedule & refuse the aisle\\
    Brake deceleration & loaded stop trial & scheduled & re-adjudicate\\
    Wear budget & hours and load cycles & counted & re-adjudicate\\\bottomrule
  \end{tabular}\\[0.2cm]
  \normalsize\uncover<2->{A boot checksum shows the calibration file is intact; only a physical check shows the camera still sits where the file says.}
  \mlsyscredit{Deployment Release, \S16.5.1 and table 1. A \val{shortfall_pct}\% weaker deceleration spends the \val{spare}~mm spare (Napkin Math 16.1).}
  \note{[2.5 min] [CORE] Definition: a standing condition binds a verdict to one premise, with a detector, a cadence set by how fast that premise degrades, and a validated response. NARRATE: read across. Fast premises are checked every \val{t_tick} ms tick; payload at each load-changing event; floor friction by survey and a spill procedure, because no sensor on the machine sees clear oil ahead of the wheels; brake wear by scheduled loaded stop trials, because a lining \val{shortfall_pct} percent weaker spends the whole spare. Calibration, the sixth row of the chapter's table, is checked by a boot self-test and by inspection as mounts wear; on failure the machine refuses until recalibrated. Facility premises need accountable people and enforceable exclusion rules as well as sensors. Step 2: the classic trap. ENGAGE: an over-the-air update preserves the calibration file bit for bit. Is calibration verified? (Expected: no; vibration and heat move the bracket; only a physical self-test checks the pose.) WARN: cadences follow physical wear; software release schedules do not set them. UL 4600 calls these watched measurements safety performance indicators, each with a threshold fixed before service. FLEX: [CORE].}
\end{frame}

\begin{frame}{Unknown refuses like known-false; their causes are logged apart}
  \centering
  \safeimg[width=0.97\textwidth,height=3.6cm,keepaspectratio]{tri.pdf}\\[0.25cm]
  {\small \texttt{is\_safe == true} reads the same for a dead sensor and a clear aisle.}
  \mlsyscredit{Deployment Release, definition 16.3 and \S16.5.2. Lecture 15: UNOBSERVABLE and INVALID trials leave gaps.}
  \note{[1.5 min] [OPTIONAL] Definition: the envelope trichotomy. Known-true: verified, fresh evidence puts the qualified state set inside the envelope, so tracking may continue. Known-false: a measured breach, such as tracking error above its \val{d_track} mm bound. Unknown: a lost encoder link, a frame failing its cyclic redundancy check (CRC), a diverged estimator; the controller cannot tell whether the machine is inside or outside. The permission decision is the same for the last two: revoke tracking and run the fallback rung Lecture 12 validated for the current state. The flight recorder keeps the causes apart, because the repairs differ. LINK: the same rule applies to test records; Lecture 15's UNOBSERVABLE and INVALID trials leave gaps in coverage. ENGAGE: which state was Challenger's joint in at \val{joint_f} F? (Expected: unknown.) WARN: a stored command is not valid forever; an active chunk runs only within its own lease and evidence. FLEX: [OPTIONAL].}
\end{frame}

\begin{frame}{Estimate: how far does the base travel blind?}
  \mlsysactivity{Estimate}{2 min}
  \centering
  \mlsyssteps[width=0.97\textwidth,height=3.7cm,keepaspectratio]{silent_step}{2}\\[0.0cm]
  \begin{columns}[T]
    \begin{column}{0.44\textwidth}
      \begin{mlsyscard}{mlsysaccent}\footnotesize\textbf{Given:} \val{v_aisle}~m/s, so \val{mm_per_ms}~mm per ms. Lease \val{t_lease}~ms (Lecture 4); tick \val{t_tick}, bus \val{t_bus}, brake onset \val{t_onset}~ms.\end{mlsyscard}
    \end{column}
    \begin{column}{0.52\textwidth}\footnotesize
      \wstep{1}{Silent proposer: \val{t_silent}~ms, \textbf{\val{d_silent}~mm}. Immediate flag: \val{t_flagged}~ms, \val{d_flagged}~mm}
      \wstep{2}{Monitoring saves up to \textbf{\val{d_saved}~mm}, but only if the flag arrives within \textbf{\val{t_crossover}~ms}}
    \end{column}
  \end{columns}
  \mlsyscredit{Deployment Release, Napkin Math 16.2 (lease and tick chosen; bus and onset illustrative; derived in calcs.py).}
  \note{[3 min] [CORE] Worked example 2. ENGAGE: two minutes; the base runs the aisle at \val{v_aisle} m/s. How far does it travel before braking when the proposer falls silent, and how far when a monitored premise fails and its flag reaches the enforcer at once? Step 1, top row: a silent proposer is caught only when its \val{t_lease} ms chunk lease (Lecture 4) runs out, so \val{t_lease} + \val{t_tick} + \val{t_bus} + \val{t_onset} = \val{t_silent} ms, \val{d_silent} mm. Middle row: a monitored premise whose flag reaches the enforcer at once takes \val{t_tick} + \val{t_bus} + \val{t_onset} = \val{t_flagged} ms, \val{d_flagged} mm. The \val{t_flagged} ms leaves out the time to detect the failure. Step 2, bottom row: a flag that takes \val{t_crossover} ms brakes no sooner than the lease. Monitoring saves up to \val{d_saved} mm only if the flag latency is under \val{t_crossover} ms, the length of the lease; a slower flag stops the base no sooner. WARN: if the case does not bound the detector's latency, it cannot claim protection from that monitor. Writing the flight-recorder log must never delay the stop. FLEX: [CORE].}
\end{frame}

\begin{frame}{A run of harmless erosions moved what reviews accepted}
  \begin{columns}[c]
    \begin{column}{0.52\textwidth}
      \centering\safeimg[width=\textwidth,height=4.3cm,keepaspectratio]{erosion.pdf}
    \end{column}
    \begin{column}{0.46\textwidth}
      \begin{logblock}[7.2]
        ``NASA and Thiokol accepted escalating risk apparently because they `got away with it last time.'\,''
      \end{logblock}
      \vspace{0.1cm}
      {\footnotesize Rogers Commission, finding 3}\\[0.3cm]
      {\small Fix each watched threshold and its response \textbf{before} service.}
    \end{column}
  \end{columns}
  \mlsyscredit{Data: Dalal, Fowlkes, and Hoadley 1989 via UCI, in launch order. Quote: Rogers Commission report, vol.~1, ch.~VI. Deployment Release, \S16.7 and \S16.9.5.}
  \note{[1.5 min] [OPTIONAL] The same \val{n_flights} flights in launch order. Distress appeared on \val{n_distressed} of them, early and late, and every flight came back, so each review treated the last success as evidence of safety. Diane Vaughan named this the normalization of deviance. On the running machine it looks like localization residuals creeping toward their allowance, age-check trips dismissed as nuisance stops, stopping distances slightly longer than the first trials, and enforcer trips cleared without a trace review. The defense: fix each safety performance indicator's threshold and response before service, so a run of harmless readings cannot move it, and repeat adversarial review on a schedule. ENGAGE: which line in the table two slides back is most likely to drift this way? (Expected: brake deceleration, checked only at scheduled trials.) FLEX: [OPTIONAL]. Land the segment: a verdict lives while every premise is known-true, checked at a cadence set by its wear.}
\end{frame}

% =============================================================================
% SECTION: UPDATES
% =============================================================================
\mlsysnote{[0.5 min] [CORE] Fleets change software often. What does the machine check before it moves, and what does an update reopen?}
\mlsysdivider{What does the machine check before it moves, and what does an update reopen?}{Deployment Release, \S16.8 and \S16.9.1 to \S16.9.4}
\section{Updates}

\begin{frame}{Waymo recalled its driving software \val{rc_words} times in \val{rc_months} months}
  \begin{columns}[c]
    \begin{column}{0.72\textwidth}
      \footnotesize
      \begin{tabular}{@{}P{1.55cm}P{1.35cm}P{5.3cm}@{}}\toprule
        \textbf{Campaign} & \textbf{Vehicles} & \textbf{NHTSA summary (excerpt)}\\\midrule
        24E-013 & \val{rc1_n} & software released before Dec.~20, 2023 ``may inaccurately predict the movement of a towed vehicle''\\
        24E-049 & \val{rc2_n} & software before May 29, 2024 ``could fail to avoid a pole or similar object''\\
        25E-034 & \val{rc3_n} & software before the Nov.~7, 2024 release ``may cause the vehicles to collide with [...] chains and gates''\\\bottomrule
      \end{tabular}\\[0.15cm]
      {\footnotesize Each remedy: ``Waymo has updated the ADS software.''}
    \end{column}
    \begin{column}{0.26\textwidth}
      \centering\safeimg[width=\textwidth,height=2.6cm,keepaspectratio]{amazon-drive-unit.jpg}\\[0.05cm]
      {\footnotesize Warehouse fleets update too: an Amazon drive unit}
    \end{column}
  \end{columns}
  \mlsyscredit{NHTSA recall records 24E013, 24E049, 25E034 (api.nhtsa.gov, retrieved 2026-10-19). Photo: Auledas, CC BY 4.0, Wikimedia Commons.}
  \note{[2 min] [CORE] Real records from NHTSA, the US vehicle-safety regulator. Each recall names a software version, so the claim and the defect attach to a configuration. Read the three rows: a towed vehicle predicted wrongly (\val{rc1_n} vehicles, filed February 2024), poles (\val{rc2_n}, June 2024), chains and gates (\val{rc3_n}, May 2025). The fleet grew \val{rc_growth} times across them (derived). In each case the remedy was a software update; for 25E-034 every vehicle was repaired by December 26, 2024, \val{rc_days} days before the report was filed (derived). ENGAGE: an update fixed the gate case. Which earlier claims did the same update reopen? (Expected: every claim whose evidence it touched: perception timing, planning, and the stops the new behavior produces.) WARN: a recall is a public, after-the-fact REFUSE; the release gate is the machine's own, before motion. The photo: warehouse fleets of drive units ship software updates too. FLEX: [CORE].}
\end{frame}

\begin{frame}{The release manifest binds every record to one configuration}
  \centering
  \mlsyssteps[width=0.98\textwidth,height=6.0cm,keepaspectratio]{manifest_step}{3}
  \mlsyscredit{After Deployment Release, \S16.8, figure 3, and table 6. Lecture numbers mark where each record was taught.}
  \note{[3 min] [CORE] The record this lecture adds. Step 1: the records of the whole course, by the lecture that taught each: limit record (2), handoff record (4), provenance (5), policy manifest (6), evaluation record (7), observation contract (8), belief record (9), intent lease (10), trajectory record (11), enforcement record (12), placement record (13), authority record (14), fault records (15). Step 2: the manifest binds reviewed snapshots of them by digest under one Merkle root (a digest is a hash that changes if any byte does; the Merkle root is one digest over all the digests), so changing any evidence leaf changes the root that is signed. It adds the root claim, the argument digest, binary hashes for policy weights, enforcer, and kernel, calibration tables, standing conditions, a wear budget, an expiry date, and the adjudicator's signature from a hardware security module. Step 3: the gate authenticates, then checks standing conditions, before motor power stages energize. Any mismatch refuses, except stale calibration covered by a separately signed COND. ENGAGE: an engineer re-flashes last month's firmware after a brake check fails. Does the gate pass? (Expected: no; old software cannot fix a worn pad.) LINK: Lecture 17's approval record will join the gate. FLEX: [CORE].}
\end{frame}

\begin{frame}{Update metadata already signs hashes and an expiry date}
  \begin{columns}[c]
    \begin{column}{0.63\textwidth}
      \begin{logblock}[8.2]
        \{\\
        \ind{1}"signatures": [ \{\\
        \ind{3}"keyid": "65171251[...]",\\
        \ind{3}"sig": "d65f8db0[...]" \} ],\\
        \ind{1}"signed": \{\\
        \ind{2}"\_type": "targets",\\
        \ind{2}\textbf{"expires": "2030-01-01T00:00:00Z"},\\
        \ind{2}"targets": \{\\
        \ind{3}"file1.txt": \{\\
        \ind{4}\textbf{"hashes": \{ "sha256": "65b8c67f[...]"}, [...] \},\\
        \ind{4}"length": 31 \}, [...]\\
        \ind{2}"version": 1 \} \}
      \end{logblock}
    \end{column}
    \begin{column}{0.35\textwidth}
      \centering\safeimg[width=\textwidth,height=2.6cm,keepaspectratio]{hsm-nshield.jpg}\\[0.05cm]
      {\footnotesize A hardware security module: the signing key never leaves it}\\[0.2cm]
      {\small A robot manifest adds calibration, standing conditions, and a wear budget.}
    \end{column}
  \end{columns}
  \mlsyscredit{The Update Framework (TUF), python-tuf tests/repository\_data/.../targets.json @ 2c771a0 (test repository; lines omitted). Photo: Alexander Klink, CC BY 3.0, Wikimedia Commons.}
  \note{[1.5 min] [OPTIONAL] A real artifact: targets metadata from The Update Framework (TUF), the design Uptane extended to vehicles, taken from python-tuf's test repository. It holds the fields a manifest needs: a signature by key ID, an expiry date, and for each file a hash and a length. A client refuses a file whose hash or length differs, and refuses expired metadata. TUF splits root, targets, and release keys across separate machines so one stolen key cannot sign firmware for every robot. The photo: a hardware security module keeps the adjudicator's private key inside its boundary. ENGAGE: what does a valid signature prove about the camera's pose? (Expected: nothing; it proves which key signed the file.) WARN: software metadata attests bytes. The robot's manifest must also bind physical evidence: calibration, standing conditions, wear. FLEX: [OPTIONAL].}
\end{frame}

\begin{frame}{An update reopens every claim whose evidence it touches}
  \centering
  \safeimg[width=0.97\textwidth,height=4.9cm,keepaspectratio]{reopen.pdf}\\[0.1cm]
  {\small An unchanged enforcer binary does not mean an unchanged containment argument.}
  \mlsyscredit{After Deployment Release, \S16.9.1, \S16.7, and \S16.8. Lecture 6: every updated checkpoint is a new policy.}
  \note{[2.5 min] [CORE] Rows are changes; columns are the evidence each reopens. Retrained policy on the same backbone: Lecture 6's rule, an updated checkpoint is a new policy, so capability evidence and contact trajectories are re-qualified. Larger backbone with the enforcer binary byte for byte identical: more memory-bus traffic and rail droop can delay the enforcer and age frames, so timing (C2), age (C4), and independence (C5) reopen. A driver or thread-priority change does the same. An OS reflash that resets calibration reopens the stop and contact claims through the localization allowance. A worn brake lining reopens the stop claim with no software change at all. ENGAGE: a team swaps the camera lens. Which columns light up? (Expected: stop and contact through calibration, and age if exposure changes.) WARN: rolling back a binary cannot repair a worn brake or a stale calibration. The image and its calibration schema are one rollback unit, and the swap happens only parked, with drives unpowered. FLEX: [CORE].}
\end{frame}

\begin{frame}{Vote: what qualifies an updated policy?}
  \mlsysactivity{Vote}{2 min}
  \normalsize True or false?\\[0.25cm]
  \choice{1}{Near-zero shadow-mode disagreement over 10,000 hours qualifies the candidate.}\uncover<2->{\enspace\textbf{\textcolor{rolered}{False}}}\\[0.2cm]
  \choice{2}{Routing 5\% of missions to the candidate, as web canaries do, bounds the risk.}\uncover<2->{\enspace\textbf{\textcolor{rolered}{False}}}\\[0.2cm]
  \choice{3}{A contained canary can pass on a non-inferiority margin fixed before the trial.}\uncover<2->{\enspace\textbf{\textcolor{rolegreen}{True}}}\\[0.25cm]
  {\footnotesize In shadow mode the candidate runs beside the incumbent with no actuator authority.}\\[0.05cm]
  \uncover<2->{\small Shadow mode filters candidates; the incumbent's actions chose every state it saw.}
  \mlsyscredit{Deployment Release, \S16.9.2 and \S16.10 (\val{m_loaded}~kg at \val{v_aisle}~m/s carries \val{ke}~J, derived). Lecture 7: shadow operation as an evidence regime.}
  \note{[3 min] [CORE] Peer instruction. Step 2: (1) false. In shadow mode the candidate runs beside the incumbent with no actuator authority and logs its action discrepancy and barrier-rate discrepancy. Every state it saw was chosen by the incumbent's commands, so agreement says nothing about the states its own commands would reach (Lecture 7's trajectory confinement). Shadow mode filters gross divergence; it cannot qualify. (2) false. Choosing fewer machines limits exposure and leaves the consequence of one action unchanged; a \val{m_loaded} kg base at \val{v_aisle} m/s carries \val{ke} J (derived), and no rollback recalls it. (3) true. A canary runs in a geofenced, derated, supervised corridor under the enforcer, and passes if the upper 95 percent bound on its intervention-rate ratio stays below 1 plus a margin declared before the trial; a sequential test can stop early. ENGAGE: why must the margin be fixed in advance? (Expected: a margin chosen after the data can make any result pass.) FLEX: [CORE]. Land the segment.}
\end{frame}

% =============================================================================
% CLOSE
% =============================================================================
\begin{frame}{Recall: what verdict should Challenger have got?}
  \mlsysactivity{Retrieval}{1.5 min}
  \begin{columns}[c]
    \begin{column}{0.30\textwidth}
      \centering\safeimg[width=\textwidth,height=1.9cm,keepaspectratio]{thumb_oring_step3.pdf}\\[0.02cm]
      {\footnotesize Flight record: \val{coldest_f} to \val{warmest_f}~°F. Launch joint: about \val{joint_f}~°F.}\\[0.2cm]
      \only<1|handout:0>{\normalsize Write Challenger's launch as a claim, a premise, and a verdict.}
    \end{column}
    \begin{column}{0.68\textwidth}\footnotesize
      \uncover<2->{\begin{tabular}{@{}P{1.9cm}P{3.3cm}P{4.0cm}@{}}\toprule
        & \textbf{Challenger} & \textbf{Mug delivery}\\\midrule
        Claim & joints seal at ignition & stop in \val{d_clear_m}~m; contact under limits\\
        Weakest premise & joint $\ge$ \val{coldest_f}~°F, the flight evidence & loaded deceleration \val{a_brake}~m/s$^2$\\
        Status & unknown at \val{joint_f}~°F & illustrative\\
        Verdict & \textbf{\textcolor{rolered}{REFUSE}} & \textbf{\textcolor{rolered}{REFUSE}}; COND at \val{v_restricted}~m/s with named trials\\
        Standing condition & joint temperature in the launch criteria & floor, payload, brake wear, calibration\\\bottomrule
      \end{tabular}}
    \end{column}
  \end{columns}
  \mlsyscredit{Deployment Release, War Story 16.1 and \S16.6; Rogers Commission report, 1986.}
  \note{[3 min] [CORE] Back to the hook. Step 1: retrieval, 90 seconds. Step 2: Challenger's claim was that the joints would seal at ignition. Its premise was a joint temperature inside the flight evidence, \val{coldest_f} F or warmer. At about \val{joint_f} F the premise was unknown, so the verdict is REFUSE, which is what Thiokol's engineers first wrote. The missing standing condition was the \val{thiokol_f} F limit; the launch director testified that had the \val{thiokol_f} F limit been in the launch commit criteria, they would not have launched. The mug case is refused today for the same structural reason, and it names the trials that would earn COND at \val{v_restricted} m/s. ENGAGE: in one sentence, why can a machine that completes every trial still be refused? (Expected: trials show what happened under conditions met; permission is a claim about conditions it will meet, and each unmeasured premise refuses.) FLEX: [CORE].}
\end{frame}

\begin{frame}{Takeaways}
  \thumbline{thumb_budget_bar_aisle_step8.pdf}{Release decides inequalities over a named envelope; the aisle stop leaves \val{spare} of \val{d_clear_mm}~mm, and \val{reserve}~mm after admission work.}\\[0.16cm]
  \thumbline{thumb_legs_step2.pdf}{Required premises multiply (\val{n_premises_words} at \val{p_premise} hold at \val{p_conj}); only legs shown to be independent add confidence.}\\[0.16cm]
  \thumbline{thumb_tree_step3.pdf}{On the book's evidence every claim refuses; a refusal names its measurement, and COND needs evidence of its own.}\\[0.16cm]
  \thumbline{thumb_silent_step2.pdf}{A verdict lives while each premise's detector reads known-true; unknown refuses like a breach.}\\[0.16cm]
  \thumbline{thumb_manifest_step3.pdf}{The manifest binds records, binaries, and calibration by digest; an update reopens what its change touches.}
  \note{[1.5 min] [CORE] One line per segment. Ask students which slide taught each line. If short on time, read the second and the fourth.}
\end{frame}

\begin{frame}{Open questions}
  \begin{columns}[c]
    \begin{column}{0.64\textwidth}
      \small
      \begin{itemize}\setlength{\itemsep}{0.12cm}
        \item A fleet ships software every week. Which claims can a weekly release re-qualify, and which need a slower cadence?
        \item Who should sign a release: the builder, an independent assessor, or a regulator? What does each signature prove?
        \item Waymo's fatal-outcome interval spans \val{fat_lo} to \val{fat_hi} percent. When is that enough evidence to expand a fleet?
      \end{itemize}
    \end{column}
    \begin{column}{0.34\textwidth}
      \centering\safeimg[width=\textwidth,height=3.2cm,keepaspectratio]{waymo-jaguar-sf.jpg}\\[0.05cm]
      {\footnotesize A driverless Waymo in San Francisco}
    \end{column}
  \end{columns}
  \mlsyscredit{Photo: Dllu, CC BY 4.0, Wikimedia Commons.}
  \note{[1 min] [OPTIONAL] No single answer. Pick one for discussion or assign one as a short written response. The first comes from modular safety cases (\S16.9.1). The second sets the adjudicator's signature against UL 4600's independent assessment and the EU Machinery Regulation's third-party assessment for learned safety functions (\S16.3.4, \S16.4). The third asks what counts as enough when rare outcomes carry the decision (\S16.9.6).}
\end{frame}

\begin{frame}{Read and explore}
  \small
  \begin{itemize}\setlength{\itemsep}{0.02cm}
    \item \textbf{Read:} Deployment Release, \S16.1 to \S16.11; Checkpoints 16.1 to 16.3; Napkin Math 16.1 and 16.2 by hand.
    \item Rogers Commission, \emph{Report on the Space Shuttle Challenger Accident}, 1986, ch.~V and VI.
    \item Kelly, \emph{Arguing Safety}, 1998: Goal Structuring Notation.
    \item Knight and Leveson, IEEE TSE 1986: independence of diverse versions.
    \item Vaughan, \emph{The Challenger Launch Decision}, 1996: normalization of deviance.
    \item Kusano et al., Traffic Injury Prevention, 2025: crash rates by outcome.
  \end{itemize}
  \note{[1 min] [CORE] The chapter first, with its checkpoints; Checkpoint 16.3 asks students to fill the measurement table for their own platform and derive a verdict. The Rogers report's chapters V and VI hold the two Thiokol documents and the flight-history figures. Kelly defines the graph notation. Knight and Leveson measured coincident failures in independently written programs. Vaughan names the drift that standing conditions guard against. Kusano et al.\ show fleet evidence reported outcome by outcome.}
\end{frame}

\begin{frame}[c]
  \begin{columns}[c]
    \begin{column}{0.34\textwidth}
      \centering\mlsyskitfig[width=\textwidth,height=4.2cm,keepaspectratio]{stack_governance}
    \end{column}
    \begin{column}{0.62\textwidth}
      {\small\textcolor{lightgray}{NEXT LECTURE \textbar{} RESPONSIBLE PHYSICAL AI}}\\[0.5cm]
      {\LARGE\bfseries Who answers for a machine that moves on its own when its motion harms someone?}
    \end{column}
  \end{columns}
  \note{[0.5 min] [CORE] The manifest names the adjudicator who accepted the configuration. But the crossing rule, the staffing of the remote assistant, the permit, and the uses of the recorded data were decided by others, and none appears in a per-tick log. Lecture 17 traces each decision to a named owner and adds the approval record that the release gate reads.}
\end{frame}

\end{document}
